\subsection{\texorpdfstring{第~\ref{sec:gauss-bonnet-theorem}~节}{Gauss–Bonnet 定理}}\label{subsec:solutions-gauss-bonnet-theorem}

\begin{description}[font=\normalfont,leftmargin=2em,labelwidth=1.5em]
\RGsolution{ex:index-continuous-under-metric-interpolation}{index-continuous-under-metric-interpolation}
显然, $\langle\cdot,\cdot\rangle_t$ 是双线性的且对称, 并满足 $\langle v,v\rangle_t\ge0$. 此外, $\langle v,v\rangle_t=0$ 当且仅当
\[
  (1-t)\lVert v\rVert_0^2+t\lVert v\rVert_1^2=0,
\]
即当且仅当 $\lVert v\rVert_0^2=\lVert v\rVert_1^2=0$, 因而 $v=0$. 为说明函数 $I_p(t)$ 连续, 考虑相对于度量 $\langle\cdot,\cdot\rangle_t$ 的正定向标准正交标架 $\{F_{1,t},F_{2,t}\}$, 使得
\[
  F_{1,t}=\frac X{\lVert X\rVert_t}.
\]
则
\[
  (\overline\omega^1)_t:=\frac{\lVert X\rVert_t}{\lVert X\rVert_1}(\overline\omega^1)_1
\]
关于 $t$ 连续. 由于 $F_{2,t}$, $\mathrm{d}(\overline\omega^1)_t$ 及 $\mathrm{d}(\overline\omega^2)_t$ 随 $t$ 连续变化, 所以
\[
  (\overline\omega_1^2)_t
  =\mathrm{d}(\overline\omega^1)_t(F_{1,t},F_{2,t})(\overline\omega^1)_t
  +\mathrm{d}(\overline\omega^2)_t(F_{1,t},F_{2,t})(\overline\omega^2)_t
\]
也随 $t$ 连续变化, 从而 $I_p(t)$ 连续.

\RGsolution{ex:hessian-and-morse-theorem}{hessian-and-morse-theorem}
\begin{enumerate}[label=(\alph*),ref=\alph*,leftmargin=*]
\item 设 $x=(x^1,\ldots,x^n)$ 是以 $p$ 为中心的局部坐标, $\widehat f:=f\circ x^{-1}$ 是 $f$ 在该局部坐标中的表达式. 由于 $p$ 是临界点, 有
\[
  \frac{\partial\widehat f}{\partial x^i}(0,\ldots,0)=0.
\]
现在,
\[
  \frac{\partial}{\partial s}(f\circ\gamma)
    =\sum_{i=1}^n\frac{\partial\widehat f}{\partial x^i}\frac{\partial\gamma^i}{\partial s}
\]
(其中 $\gamma^i:=x^i\circ\gamma$), 因而
\[
  \frac{\partial^2}{\partial t\partial s}(f\circ\gamma)
  =\sum_{i,j=1}^n\frac{\partial^2\widehat f}{\partial x^i\partial x^j}
     \frac{\partial\gamma^i}{\partial s}\frac{\partial\gamma^j}{\partial t}
  +\sum_{i=1}^n\frac{\partial\widehat f}{\partial x^i}\frac{\partial^2\gamma^i}{\partial t\partial s}.
\]
取 $(s,t)=(0,0)$, 得
\[
  (Hf)_p(v,w)=\sum_{i,j=1}^n\frac{\partial^2\widehat f}{\partial x^i\partial x^j}(0,\ldots,0)v^iw^j,
\]
这说明 $(Hf)_p$ 的确是一个对称 2-张量. 由于它只依赖于 $v$ 和 $w$ 的分量, 所以它是良定义的 (即与映射 $\gamma$ 的选取无关).
\item 在以某个临界点 $p\in M$ 为中心的局部坐标 $x=(x^1,\ldots,x^n)$ 中, 有
\[
  \widehat f(x^1,\ldots,x^n)=f(p)+\sum_{i,j=1}^n\frac12H_{ij}x^ix^j+o(\lVert x\rVert^2),
\]
其中 $\widehat f:=f\circ x^{-1}$, 而
\[
  H_{ij}:=\frac{\partial^2\widehat f}{\partial x^i\partial x^j}(0,\ldots,0)
\]
是 Hessian 的分量, 因而
\[
  \frac{\partial\widehat f}{\partial x^i}(x^1,\ldots,x^n)
    =\sum_{j=1}^nH_{ij}x^j+o(\lVert x\rVert).
\]
由于 Hessian 非退化, 矩阵 $(H_{ij})$ 可逆, 因此存在 $p$ 的一个邻域, 在其中 $p$ 是 $f$ 的唯一临界点. 所以 $f$ 的临界点是孤立的, 并且由于 $M$ 紧, 这样的点只能有有限多个 (否则它们会在一个非孤立临界点处聚集). 由于 $(H_{ij})$ 对称且非退化, 存在线性坐标变换把它化为形如 $\operatorname{diag}(1,\ldots,1,-1,\ldots,-1)$ 的对角矩阵. 若 $M$ 是 2 维的, 可能性为 $\operatorname{diag}(1,1)$ (此时 $p$ 为极小点), $\operatorname{diag}(-1,-1)$ (此时 $p$ 为极大点) 及 $\operatorname{diag}(1,-1)$ (此时 $p$ 为鞍点). 在 $M$ 上选一个 Riemann 度量, 使其在每个临界点的某邻域内与这些坐标下的 Euclid 度量一致, 则在 $p$ 附近, 若 $p$ 为极小点, 有
\[
  \operatorname{grad} f=(x^1+o(\lVert x\rVert))\frac{\partial}{\partial x^1}
                       +(x^2+o(\lVert x\rVert))\frac{\partial}{\partial x^2};
\]
若 $p$ 为极大点, 有
\[
  \operatorname{grad} f=(-x^1+o(\lVert x\rVert))\frac{\partial}{\partial x^1}
                       +(-x^2+o(\lVert x\rVert))\frac{\partial}{\partial x^2};
\]
若 $p$ 为鞍点, 有
\[
  \operatorname{grad} f=(x^1+o(\lVert x\rVert))\frac{\partial}{\partial x^1}
                       +(-x^2+o(\lVert x\rVert))\frac{\partial}{\partial x^2}.
\]
由 Example~\ref{ex:four-vector-field-indices} 可知, 当 $p$ 为极大点或极小点时, $\operatorname{grad} f$ 在 $p$ 处的指数为 1, 当 $p$ 为鞍点时指数为 $-1$. 因为 $\operatorname{grad} f$ 的零点恰为 $f$ 的临界点, 所以
\[
  \chi(M)=m-s+n.
\]
\end{enumerate}

\RGsolution{ex:gauss-curvature-by-geodesic-triangle}{gauss-curvature-by-geodesic-triangle}
注意, 尽管 $\partial\Delta$ 不是光滑流形, 仍可通过\textquotedblleft 磨圆角点\textquotedblright 用一列光滑流形逼近它, 然后取极限. 因此下面将把 $\partial\Delta$ 当作光滑流形处理.
\begin{enumerate}[label=(\alph*),ref=\alph*,leftmargin=*]
\item 考虑沿 $\partial\Delta$ 平行的向量场 $V$. 因为 $\Delta$ 的边是测地线, $\partial\Delta$ 的切向量仅在顶点处相对于 $V$ 旋转, 旋转总量为
\[
  (\pi-\alpha)+(\pi-\beta)+(\pi-\gamma).
\]
这必等于 $2\pi$ 减去 $V$ 旋转的角度. 因此
\[
  3\pi-\alpha-\beta-\gamma=2\pi-\int_\Delta K.
\]
\item 刚才已看到, $\partial\Delta$ 的切向量相对于平行向量 $V$ 旋转的总量为
\[
  \int_{\partial\Delta} k_g=(\pi-\alpha)+(\pi-\beta)+(\pi-\gamma).
\]
由带边界流形的 Gauss–Bonnet 定理, 有
\[
  3\pi-\alpha-\beta-\gamma+\int_\Delta K=2\pi\chi(\Delta).
\]
因为 $\Delta$ 与圆盘同胚, $\chi(\Delta)=1$, 所求结论得证.
\end{enumerate}
\end{description}
